\section{Related Work}
\label{sec:related}

\para{Glitch token detection.}
The phenomenon of anomalous tokens in language models was first widely publicized by \citet{rumbelow2023solidgoldmagikarp}, who demonstrated that specific tokens such as \magikarp caused ChatGPT to produce evasive, hallucinatory, or hostile responses. This observation motivated systematic detection methods. \citet{land2024fishing} proposed the first automated approach, using embedding-based indicators---specifically, near-zero $L_2$ norms from weight decay---to identify under-trained tokens across 88 open-weight models. They found that 0.1--1\% of vocabulary entries in every tested model were under-trained. \citet{li2024glitch} introduced GlitchHunter, a clustering-based approach that also contributed a taxonomy of five token types and five failure behaviors, achieving 99.4\% precision. \citet{zhang2024glitchprober} used PCA on attention patterns combined with SVM classification, while \citet{wu2024glitchminer} proposed gradient-based discrete optimization to find tokens maximizing output entropy. These methods all focus on \emph{detecting} problematic single tokens in open-weight models. We complement this line of work by testing whether detected \glitchtokens still cause failures in the latest closed-source models, and by extending the evaluation from single tokens to full passages.

\para{Token-level vulnerabilities and adversarial inputs.}
Beyond individual \glitchtokens, recent work has shown that \emph{combinations} of tokens can create failures. \citet{jang2024improbable} demonstrated that ``improbable bigrams''---pairs of tokens that are individually well-trained but rarely co-occur---cause hallucination rates of 33--77\% in models with byte-level tokenizers, compared to 0--20\% for normal text. \citet{yang2024problematic} showed that GPT-4o's tokenizer creates systematic failures for Chinese and Korean text due to undertrained tokens arising from misaligned vocabulary construction. More broadly, \citet{geiping2024coercing} systematized adversarial attacks on LLMs, identifying \glitchtokens as an active attack vector for jailbreaking and denial-of-service. Our work extends beyond tokenizer-level vulnerabilities to study how text properties such as length, content type, and error density affect reproduction fidelity.

\para{Robustness to input perturbations.}
A related body of work studies how models respond to perturbed inputs. \citet{moradi2021evaluating} showed that BERT, XLNet, and RoBERTa are sensitive to character-level perturbations such as typos and misspellings, with small changes causing significant performance drops on downstream tasks. However, this work evaluates \emph{task performance} under perturbation, not \emph{reproduction fidelity}---whether a model can copy perturbed text as-is. Our misspelling experiments directly test this distinction, asking whether models preserve deliberate errors or silently correct them.

\para{Generalization failures in language models.}
More broadly, our work connects to research on systematic generalization failures. \citet{berglund2023reversal} demonstrated the ``reversal curse,'' where models trained on ``A is B'' fail to learn ``B is A.'' While this is a different type of failure than reproduction errors, both reveal fundamental asymmetries in how models process information. Our finding that \gptfull reproduces misspelled text but truncates clean text at the same length represents an analogous asymmetry in the reproduction domain.
